
Formal verification feedback loops have been shown to improve LLM performance on theorem proving and code verification tasks. The underlying mechanism remains unresolved: does structured formal feedback function as a \textit{local logical oracle} that directs probability mass toward premise-consistent tactics, or as a \textit{policy regularizer} that sharpens the model regardless of the semantic content of the feedback signal? This paper introduces the oracle/regularizer distinction as a testable mechanistic question, proposes the \textbf{locality score} — the fraction of post-DPO probability mass shift concentrated on premise-consistent tactic categories — as a direct mechanistic probe, and describes a 3-condition experiment (grounded, ungrounded, permuted control) within the BFS-Prover framework designed to test this distinction. An attempt to execute this experiment encountered an infrastructure failure: the LeanDojo proof state extraction layer silently substituted synthetic hardcoded error strings for real Lean4 compiler outputs, rendering all locality score measurements uninformative. Locality scores for all three conditions were uniformly 0.0000 on both the miniF2F and Vericoding evaluation datasets, with a one-sided t-statistic of 0.0000 and p-value of 1.0000. This outcome is a synthetic artifact rather than a scientific finding: the DPO training infrastructure (β=10, 100\% state alignment on synthetic data, correct loss implementation) was validated, but no real Lean4 compiler invocations occurred. The experimental design and implementation are ready for re-execution after the LeanDojo environment is properly installed. This paper documents the silent synthetic fallback failure mode and proposes pre-run environment validation gates as a methodological recommendation for LLM experiments that depend on external formal tools.
